\section{First derivative across a sum-map split}
\label{sec:h-jet}

Lemma~\ref{lem:mu} is the statement.
The first derivative of the summed amplitude jumps by
\(i\bigl(s+\sum_C\psi\bigr)\) for every value of the fibre coordinate \(d\).
The cover enters only through the common-neighbour set \(C\).
The second derivative sees \(d\) with coefficient
\(-\tfrac12(|X_i|-|X_j|)\).
A uniform nonzero vacuum fails the first-order condition on every cover.
A stored replay satisfies every order because it is the predecessor.
You could treat \eqref{eq:mudot} as a necessary condition on any jet match
that is meant to be \(C^1\) for the summed amplitude.
It is two real constraints, hence the codimension discussed in
Section~\ref{sec:codim}, and it is the same \(s\) that appears in \(P_-\).
Corollary~\ref{cor:balanced} is the symmetric case:
equal exclusive cardinalities remove \(d\) from the second derivative,
so the circle of Proposition~\ref{prop:circle} is still free at that order.
The current on the circle is Lemma~\ref{lem:BJ}.
The covers that enter \(C\) are the orbits of Proposition~\ref{prop:covers}.
How many of them have equal exclusive cardinalities, and therefore drop
the fibre coordinate out of \(\ddot\mu(t_0)\), is
Proposition~\ref{prop:balanced-count}.
Their proportion tends to zero with the degree.
The mean square of the imbalance \(\delta=|X_i|-|X_j|\) on the same
measure is Corollary~\ref{cor:imbalance}.
